\documentclass[11pt,a4paper]{article}
\usepackage[margin=22mm]{geometry}
\usepackage{amsmath,amssymb,booktabs,array}
\usepackage[hidelinks]{hyperref}
\setlength{\parindent}{0pt}\setlength{\parskip}{5pt}
\raggedright\emergencystretch=3em
\begin{document}
\begin{center}{\Large P.A priorities, a joint field certificate and donor-side completion}\\[4pt]
8 October 2026 --- methodology-only analytical review\end{center}
\textbf{Status.} P.A is approved. T.A amount ceilings stay fixed. D1/D2 below
are conditional donor completions; neither is adopted. No whole solver,
clearance or gradient result is claimed.

\section*{1. Approved meaning and static rate problem}
Use metres, seconds and continuous vehicle-equivalent mass (eq). The common
triangular flux is
\begin{equation}
 F(\rho)=\min\{u\rho,w(\kappa-\rho)\},\quad u=30,\ w=5,\ \kappa=3/20,
 \quad\rho_c=3/140,\ Q=9/14,\quad L=30,\ T=1.
\end{equation}
At the regular clock donor $a$ has sampled sending mass $S_a$ (eq), choice
fractions $\theta_{ab}$, and desired amounts/weights
\begin{equation}
 d_{ab}=\mu_{ab}=\theta_{ab}S_a,\quad \sum_b\theta_{ab}=1,\quad
 b_{ab}(t)=d_{ab}-\int_0^t f_{ab}(s)ds\ge0.
\end{equation}
P.A freezes $\mu$ within this public interval; queue growth does not increase
its priority. Weights refresh at the original next request clock. Ramp share
$p$ remains geometric/fixed. Active ports have $b_{ab}>0$ and original
$\mu_{ab}>0$; fulfilled ports have flow zero until refresh. No mass belongs
to a request or effect record.

For one instantaneous own-field state use physical sending demands $D_a$ and
capacities $Q_a$ (eq/s), and own target receiving $R_b$ (eq/s). Joint bounds are
\begin{equation}
 f_{ab}\ge0,\quad\sum_b f_{ab}\le D_a\le Q_a,\quad
 \sum_a f_{ab}\le R_b.
\end{equation}
Independent target awards cannot spend the same donor twice. Two candidate
port problems differ in oriented potential ceilings:
\begin{equation}
 \text{D1: } f_{ab}\le\theta_{ab}Q_a;\qquad
 \text{D2: no additional oriented ceiling (the shared donor bound applies).}
\end{equation}
These are not the held nominal bounds $f_{ab}\le d_{ab}/T$ rejected at T.
D1 has source-disaggregated lineage but is an explicit richer-field rate
interpretation; P.A itself did not approve it.

For the ramp target A, let $D_r^+$ be ramp physical sending if its claim is
active, else zero. N.A permits a convenient algebraic mainline cap
\begin{equation}
 C_A=\max\{(1-p)R_A,R_A-D_r^+\}.
\end{equation}
Allocate mainline into A up to $C_A$, with fixed P.A weights; if its actual
inflow is $m\le C_A$, ramp entry is $f_r=\min\{D_r^+,R_A-m\}$. This is one
shared $R_A$, not two receiving resources. It implements fixed priority and
unused-share transfer. A strictly restricted ramp has $D_r^+>pR_A$
(its accepted flow can exceed $pR_A$ when mainline leaves unused space).
Thus $C_A=(1-p)R_A$; increasing ramp sending to capacity leaves it unchanged.

\newpage
\section*{2. Finite conditional port construction and amount events}
For either declared D problem, start all active rates at zero. Increase
still-free port rates in their fixed proportions: $df_{ab}=\mu_{ab}\,d\lambda$.
The water level $\lambda$ has units s$^{-1}$ and is \emph{not physical time}.
For the current free set $E$, the next increment is the minimum of
\begin{align}
 \delta_a&=\frac{D_a-\sum_b f_{ab}}{\sum_{(a,b)\in E}\mu_{ab}},\qquad
 \delta_b=\frac{R_b-\sum_a f_{ab}}{\sum_{(a,b)\in E}\mu_{ab}},\\
 \delta_{ab}&=\frac{\theta_{ab}Q_a-f_{ab}}{\mu_{ab}}
 \quad\text{(D1 only)}.
\end{align}
Use $C_A$ in place of $R_A$ for the mainline subproblem. Omit resources with
no free incident port. After the increment, freeze every free port incident
to any tied saturated donor/recipient; under D1 also freeze every tied
oriented ceiling. Resume the other ports. At least one port freezes each
pass, so at most $|E_0|$ passes are needed. Ties are simultaneous, without
lexicographic traffic priority or perturbing input values.

This constructs the unique capped proportional sharing result for its
\emph{declared} donor/recipient sharing conventions: before the earliest
binding level, every active port must grow in its prescribed ratio. That
level is forced by the first resource bound. Its incident rates are fixed;
repeat on residual resources. This finite induction specifies all rates.
It is not a theorem that P.A alone chooses the directional ceiling, nor an
imported generic FIFO/network-throughput optimum.

For any strictly restricted donor ($\sum_bf_{ab}<D_a$), its physical row bound
never froze a port. Raising that unused bound to $Q_a$ cannot change the
result. This is the required local invariance after its own queue forms.
Recipient capacity bounds never frozen are likewise irrelevant when raised.
Under D1, orientation caps remain explicit independent restrictions.

A cap/weight can govern static rate allocation without owning traffic.
Actual rates integrate in physical time until the next field or allowance
change. For locally constant positive rates, an allowance exhaustion is at
\begin{equation}
 \tau_{\rm amount}=\min_{ab:f_{ab}>0}\frac{b_{ab}}{f_{ab}}.
\end{equation}
At simultaneous exhaustion remove every fulfilled port; retain other original
weights and amounts. Do not normalize $\theta$, enlarge $d$, resample choice,
replace $b$ by a nominal rate, or create effects early. Re-solve with the
own field. This static construction is conditional on D; all-mode wave
continuation and numerical/reference adoption remain unqualified.

\newpage
\section*{3. A joint local one-step field and transfer certificate}
Geometry: donor1 through to A; donor2 lateral to A and through to B; finite
ramp to A; B through to next section C. Meter/ramp and donor2 inlets are
closed, outside demand/queue zero. Donor1 inlet is $1/5$. C is initially
empty. Original source masses are
\begin{equation}
 S_1=1/5,\quad S_2=1/10,\quad S_r=S_B=Q,\quad
 \theta_{2A}=1/4,\ \theta_{2B}=3/4.
\end{equation}
All other enabled diagnostic choices are through. Supplied prior Edie
observations give these choices; no reached history is claimed.

A starts at density $r_A=9/100$ on $[0,1]$ and vacuum ahead, with an old
record desired speed $s_A=10/3$. B starts at $r_B=11/100$ on $[0,29]$ and
vacuum ahead, old speed $s_B=20/11$. Their zero-relative-passing own queues
have $F(r_A)=s_A r_A=3/10$, $F(r_B)=s_Br_B=1/5$. Use illustrative $p=1/3$:
\begin{equation}
 f_r=1/10,\quad C_A=1/5,\quad\mu_{1A}:\mu_{2A}=1/5:1/40=8:1,
 \quad f_{1A}=8/45,\ f_{2A}=1/45.
\end{equation}
Through B receives $1/5$ until its donor2 allowance $d_{2B}=3/40$ is used at
$t_b=3/8$. Afterwards $f_{2B}=0$. D1's oriented ceilings are nonbinding in
this specimen; D1/D2 return identical rates. A receiving stays $3/10$.
B's upstream receiving rises to $Q$ after its new vacuum rear leaves the
inlet at $t_b$, but its fulfilled donor2 request cannot restart.

Donor1 initially has uniform free density $1/150$. Its own congested exit
trace is $103/900$, shock speed $-20/97$. At $t=1$ its density integral is
$2/9=1/5+(1/5-8/45)$. The shock never reaches its inlet.

Donor2 is initially jammed on $[88/3,30]$, mass $1/10$. Before $t_b$, its
total exit flow is $g_0=2/9$ and exit trace $r_0=19/180$, rear speed
$c_0=40/19$. The release contact $-5$ reaches its initial rear at $t_0=2/15$.
At $t_b$, that rear is $567/19$. Through exhaustion leaves only lateral flow
$g_1=1/45$, with trace $r_1=131/900$. The new $-5$ contact catches the rear at
\begin{equation}
 t_c=143/360,\quad x_c=269/9,\quad c_1=g_1/r_1=20/131.
\end{equation}
The completed rear is $7855/262$ and donor2 mass is
\begin{equation}
 r_1(30-7855/262)=1/360=1/10-3/40-1/45.
\end{equation}
The rear does not reach the exit until $9/8>1$, so physical sending remains
available during the whole public interval. The resulting waves determine
that availability, rather than an external flow oracle.

The ramp initially has critical density throughout. Its closed inlet sends
a vacuum front at $u$, meeting the $-w$ exit queue front at $6/7$ s. The exit
trace is $13/100$ and final vacuum rear is $2350/91$, giving
\begin{equation}
 m_r(1)=\tfrac{13}{100}(30-2350/91)=19/35=Q-1/10.
\end{equation}

\newpage
\section*{4. Actual crossing, own counts and completed E1}
A's field is $r_A$ on $[0,1+s_A t]$ and vacuum ahead. Consequently
\begin{equation}
 m_A(t)=9/100+(3/10)t=C_A^{\rm count}(t),\quad
 N_A(t,z_A(t))=0.
\end{equation}
The superscript distinguishes this incoming count gauge from section3's
mainline receiving cap. The old record sees vacuum ahead, so K.A permits its
held speed with no extra slowdown.

B's vacuum rear after through-input exhaustion is
$b_B(t)=s_B\max\{0,t-3/8\}$. Its old record crosses the section edge at
$t_*=11/20$, then moves in C at $z_C=s_B(t-t_*)$. The explicit fields are
\begin{equation}
 \rho_B=r_B\text{ on }[b_B(t),\min(29+s_Bt,30)],\quad
 \rho_C=r_B\text{ on }[0,s_B\max\{0,t-t_*\}],
\end{equation}
with vacuum elsewhere. At crossing, the C record initially occupies zero
length; its own queue gives C receiving $1/5$, paired B-to-C flow $1/5$ for
$t>t_*$. No traffic bypass or second receiving resource appears.
\begin{align}
 I_B(t)&=\tfrac15\min(t,3/8),\quad O_B(t)=\tfrac15\max(0,t-11/20),\\
 m_B(t)&=319/100+I_B(t)-O_B(t),\qquad m_C(t)=O_B(t).
\end{align}
Before crossing $N_B(t,z_B)=0$; at crossing rebase to $N_C(t_*,0)=0$;
afterwards $N_C(t,z_C)=0$. These equal values arise from the chosen gauges;
the general rebase is $h^{\rm new}=N_{\rm receiver}(t_*,0)$, not mass transfer.

The six completed inventories are
\begin{center}\begin{tabular}{cccccc}\toprule
Donor1&Donor2&Ramp&A&B&C\\\midrule
$2/9$&$1/360$&$19/35$&$39/100$&$127/40$&$9/100$\\\bottomrule
\end{tabular}\end{center}
Their total is $774/175=739/175+1/5$: initial plus the only external input.
Every port debit has one paired credit; no exterior exit occurs by this clock.
The one-interval population integral is $1513/350$ eq s. This accounting is
not a control-benefit witness, since total inventory is unaffected by internal
redistribution when the diagnostic has no external exit.

Illustrative E1 counters give
\begin{equation}
 H_r(1)=47/50+1/10=26/25,\quad H_L(1)=49/50+1/45=451/450.
\end{equation}
With unit descriptor threshold, each creates one effect, carrying remainders
$1/25$ and $1/450$. Both attach at A's entrance with completed label $39/100$,
but retain distinct request-phase speeds (e.g. 10 and 15) and signatures.
They own zero population. A's ahead speed $F(r_A)/r_A=10/3$ caps both initial
actual speeds at the next interval; same count/position creates no invented
vehicle order. Old ending speeds are stored separately and positive recovery
uses the next regular source recurrence and its approved grade convention.

This is a completed \emph{local} sourced field/node/motion/crossing/reset
certificate, not all encounters, indefinite continuation, reached policy,
clearance or a selected A1/A2/C numerical realization.

\newpage
\section*{5. Early-depletion diagnostic: same allowed mass, different timing}
One donor starts jammed on $[86/3,30]$, mass $1/5$, inlet closed. Through
receiving is zero; lateral recipient empty/free. Original choice fraction
$1/4$ gives lateral allowance $d_L=1/20$ and through $3/20$.

D1 uses lateral $\beta_1=Q/4=9/56$ until the allowance is fulfilled. Its
trace is $33/280$, rear speed $15/11$, release contact time $4/15$, and
\begin{equation}
 t_{q1}=14/45,\quad x_{q1}=316/11,\quad
 t_{c1}=23/45,\quad x_{c1}=29.
\end{equation}
After the stop, a jam contact catches the moving rear; final donor is jammed
on $[29,30]$ with mass $3/20$. This reuses the preserved earlier certificate.

D2 uses full physical sending $\beta_2=Q$ until the same allowance is fulfilled:
\begin{equation}
 t_{q2}=7/90<4/15,\quad r_2=\rho_c,
 \quad t_{c2}=4/15+1/90=5/18,\quad x_{c2}=29.
\end{equation}
Before $t_{q2}$ the release contact is at $30-5t$. The new stop contact is
at $30-5(t-t_{q2})$: both travel at $-5$, with a critical-density slug
between. At $4/15$ the release contact reaches the initial rear; that vacuum
rear travels at $u$ until meeting the stop contact $1/90$ later. Thereafter
the jam rear is stationary at 29. This is an admissible triangular-KW field
with exactly $1/20$ departed; no retry/rerouting is needed.

At $T=1$ the empty lateral recipient has respectively
\begin{equation}
 \text{D1: }\rho=3/560\text{ on }[62/3,30];\qquad
 \text{D2: }\rho=3/140\text{ on }[83/3,30].
\end{equation}
Both masses are $1/20$ and neither has positive-measure exit before the clock.
The profiles and event times differ despite identical ending donor/recipient
mass. The proper start/stop post-event boundary trace has zero spatial
measure at the event; it changes no density integral.

This does not select a controller or demonstrate hidden information under
a shared law. It proves that the instantaneous sender completion matters.
T.A and P.A can hold in both cases.

\newpage
\section*{6. Late arrival changes accepted mass and next effects}
Initial donor is jammed on $[0,3]$ and empty ahead, inlet closed. The original
sending mass is $S=9/20$. Choice fraction $1/4$ sets
\begin{equation}
 d_L=9/80,\quad d_T=27/80.
\end{equation}
Through receiver is jammed with zero receiving; lateral receiver empty/free.
The donor's free front reaches its exit at $t_a=9/10$, and its inlet release
reaches zero at $3/5$. The free rear then travels at $u$ and is at 12 at
$T=1$. There is no source or early arrival invented by the node.

The two rate histories over the last $1/10$ second give
\begin{equation}
 M_{L1}=(Q/4)(1/10)=9/560,\qquad
 M_{L2}=Q(1/10)=9/140,
 \qquad M_{Lk}<d_L.
\end{equation}
D1 creates a congested exit trace $33/280$; its backward contact is at $59/2$
at the clock. The exact donor mass is
\begin{equation}
 \rho_c(59/2-12)+(33/280)(30-59/2)=243/560
 =9/20-9/560.
\end{equation}
D2 has critical density on $[12,30]$, mass $27/70=9/20-9/140$.
Recipient fields are $3/560$ and $3/140$ on $[0,3]$, respectively;
their integrals equal the paired lateral entries. Both are capacity/storage
admissible and respect original mass allowances.

With old eligible counter $H_0=47/50$,
\begin{equation}
 H_1=2677/2800<1,\qquad H_2=703/700>1.
\end{equation}
Only D2 creates an E1 effect at the next regular update. Birth timing/rule
itself is unchanged; this is donor-side transfer non-equivalence, not a
same-law hidden-state or gradient-usefulness result.

\section*{Decision and evidence limits}
Q.J.F.D is a shared source/availability contract, distinct from closed P.A
receiving competition and T.A integrated ceilings. D2 is provisionally
recommended to avoid an extra instantaneous directional restriction without
an established physical requirement. D1 has the source sampled-rate lineage
and remains viable. Neither is adopted. After D, qualify selected remaining
wave/contact/crossing/repeated-reset modes; full Q/J.F and Stage3 adoption
remain pending. The finite static allocation proof is not whole-network
PDE/record existence or numerical convergence. Stage4 stays unopened.

Primary source: Laval--Daganzo, eqs7--13; Wright et al.
\href{https://arxiv.org/pdf/1601.01054}{priority/node manuscript} and
\href{https://arxiv.org/pdf/1608.07623}{node execution manuscript}; Jin
\href{https://arxiv.org/pdf/1204.6727}{stationary/interior wave framework}.
Their generic FIFO/optimization theorems are not transplanted. Sections2--6
contain reviewer derivations under explicit local/conditional contracts.
\end{document}
